% 分野別類題: 変数分離形 解答例
% コンパイル: TEXINPUTS="../../../00_common:$TEXINPUTS" latexmk -xelatex answers.tex
\documentclass[a4paper,11pt]{article}
\usepackage{preamble}
\usepackage{shoakubox}

\begin{document}

\kamokumei{類題演習 解答例 --- 変数分離形}

\daimon{【解法】変数分離形の解き方と導出}

\noindent
変数分離形とは、次のように右辺が $x$ のみの関数と $y$ のみの関数の積に分解できる1階微分方程式である。
\[
\frac{dy}{dx} = f(x)\,g(y)
\]
$g(y) \neq 0$ の範囲では、両辺を $g(y)$ で割って $dx$ を掛けることで、$x$ の項と $y$ の項を左右に分離できる。
\[
\frac{dy}{g(y)} = f(x)\,dx
\]
両辺をそれぞれの変数について不定積分すると
\[
\int \frac{dy}{g(y)} = \int f(x)\,dx + C
\]
となり、これが一般解（陰関数形）を与える。可能であれば $y$ について解いて陽関数形に直す。なお $g(y) = 0$ となる $y = y_{0}$（定数）も方程式の解（特異解）になっている場合があるため、その扱いにも注意する。

また、$\dfrac{dy}{dx} = F(x+y)$ のように右辺が $x+y$（あるいは $ax+by+c$ 等）のみの関数である場合は、$u = x+y$ とおくと
\[
\frac{du}{dx} = 1 + \frac{dy}{dx} = 1 + F(u)
\]
となり、$u$ と $x$ の変数分離形に帰着できる。解いた後で $y = u - x$ に戻せばよい。初期条件が与えられている場合は、一般解に代入して任意定数を定める。

\clearpage
\begin{mondaiboxS}{問1}
$\dfrac{dy}{dx} = xy^{2}$ の一般解を求めよ。
\end{mondaiboxS}
\kaitou
$y = 0$ は明らかに解である（特異解）。$y \neq 0$ のとき変数分離して
\[
\frac{dy}{y^{2}} = x\,dx
\]
両辺を積分すると
\[
-\frac{1}{y} = \frac{x^{2}}{2} + C_{1}
\quad\Longrightarrow\quad
\frac{1}{y} = -\frac{x^{2}}{2} - C_{1}
\]
任意定数を $C = -2C_{1}$ と置き直すと
\[
\boxed{\; y = \frac{-2}{x^{2}+C} \;}
\qquad (C \text{ は任意定数。ほかに特異解 } y=0)
\]
（検算: $y=\dfrac{-2}{x^{2}+C}$ より $y' = \dfrac{4x}{(x^{2}+C)^{2}}$。一方 $xy^{2} = x\cdot\dfrac{4}{(x^{2}+C)^{2}} = \dfrac{4x}{(x^{2}+C)^{2}}$ となり両者は一致する。）

\clearpage
\begin{mondaiboxS}{問2}
$\dfrac{dy}{dx} = \dfrac{1+y^{2}}{1+x^{2}}$ の一般解を求めよ。
\end{mondaiboxS}
\kaitou
変数分離して
\[
\frac{dy}{1+y^{2}} = \frac{dx}{1+x^{2}}
\]
両辺を積分すると
\[
\arctan y = \arctan x + C
\]
よって
\[
\boxed{\; y = \tan(\arctan x + C) \;}
\qquad (C \text{ は任意定数})
\]
（検算: $\arctan y = \arctan x + C$ の両辺を $x$ で微分すると $\dfrac{y'}{1+y^{2}} = \dfrac{1}{1+x^{2}}$、すなわち $y' = \dfrac{1+y^{2}}{1+x^{2}}$ となり元の方程式に一致する。）

\clearpage
\begin{mondaiboxS}{問3}
$\dfrac{dy}{dx} = (x+y)^{2}, \quad y(0) = 0$ を解け。
\end{mondaiboxS}
\kaitou
$u = x + y$ とおくと $\dfrac{du}{dx} = 1 + \dfrac{dy}{dx} = 1 + u^{2}$ となり、変数分離形に帰着する。
\[
\frac{du}{1+u^{2}} = dx
\quad\Longrightarrow\quad
\arctan u = x + C_{1}
\quad\Longrightarrow\quad
u = \tan(x+C_{1})
\]
$u = x+y$ に戻すと $y = \tan(x+C_{1}) - x$。初期条件 $y(0) = 0$ より $\tan C_{1} = 0$ であるから $C_{1} = 0$ ととれる。したがって
\[
\boxed{\; y = \tan x - x \;}
\]
（検算: $y' = \sec^{2}x - 1 = \tan^{2}x$ であり、$(x+y)^{2} = (\tan x)^{2} = \tan^{2}x$ となって一致する。また $y(0) = \tan 0 - 0 = 0$ で初期条件も満たす。）

\end{document}
